% !TeX root = ../../main.tex
\subsubsection{CRISPulator FACS benchmark}

We used CRISPulator 0.5.1's documented interfaces for library construction,
transfection, FACS selection, and sequencing \citep{nagy2017crispulator}.
The committed Julia environment pins the tagged source revision.  For each
simulation seed, one CRISPRn guide library was constructed and reused across
four independently seeded transfection, sorting, and sequencing replicates.
The default CRISPulator phenotype mixture was retained: 75\% inactive, 5\%
negative-control, 10\% phenotype-increasing, and 10\%
phenotype-decreasing genes.  The headline screens contained 10,000 genes and
five guides per gene, so 50,000 guides, which is the order of a genome-wide
library rather than a pilot.  The default guide-quality mixture---90\%
complete-knockout guides and 10\% low-activity guides---was retained.
Transfection representation was 500 cells per guide, Gaussian phenotype noise
was 2, sorting representation was 50 cells per guide, and sequencing depth was
50 reads per guide per sample.  Multiplicity of infection was 0.20 for the
main analysis and 0.30 for the supplementary analysis; 0.20 is the point at
which the single-integration approximation these guide-level models share is
most defensible, and 0.30 stresses it.  Both multiplicity of infection and the
high-quality-guide fraction are exposed as simulation parameters, as is the
number of genes and the number of independent screen replicates.
The three genome-scale benchmark seeds were 20250724 through 20250726; the
supporting 400-gene analyses below use 20250724 through 20250728.

Two BARCS gene statistics are carried through the genome-scale ablation.
BARCS-ST is the historical calibrated signed-$z$ statistic.  BARCS-MOD reuses
the identical guide fits after guide-dispersion moderation, so the two differ
only in how the guide variance is estimated, not in the gene combiner.  The
exchangeable-normal, random-effects, and moderated random-effects statistics
are retained for the real-data comparisons but are not carried through the
genome-scale simulation, where the question is whether moderation changes the
calibration--power tradeoff rather than which of several pooling rules is
best.  Official MAGeCK-MLE 0.5.9.5 and CRISPhieRmix are included in that run,
but this restricted set is not treated as an exhaustive comparator panel.
MAGeCK-MLE's ordered marker is affinely mapped onto zero--one so that the low
samples of the reference
replicate form the all-zero reference row the MAGeCK initializer requires.
Because MAGeCK-MLE omits negative-control genes from its gene summary once
their sgRNAs are declared, all methods in the genome-scale benchmark are
scored on the genes every method returned a finite result for, 9,475--9,517
of 10,000 per run.  The negative-control diagnostic is instead taken from
each method's own output, and is therefore unavailable for MAGeCK-MLE.

We additionally performed a one-factor-at-a-time sensitivity analysis around
the baseline $(\mathrm{MOI},q,G,R)=(0.25,0.90,400,4)$, where $q$ is the
high-quality-guide fraction, $G$ is the number of genes, and $R$ is the number
of independent screen replicates.  The evaluated levels were
$\mathrm{MOI}\in\{0.10,0.25,0.40\}$,
$q\in\{0.60,0.75,0.90\}$, $G\in\{100,400,1000\}$, and
$R\in\{1,3,4,6\}$.  After removing duplicate baseline settings, this yielded
ten scenarios; each used the same five prespecified seeds, for 50 simulations.
One parameter changed at a time, so its main sensitivity remained
interpretable.  For each seed and scenario, we paired BARCS and MAGeCK-MLE
results from the low--bulk--high design and reported differences as BARCS
minus MAGeCK-MLE.  The executable benchmark also supports the complete
$3\times3\times3\times4$ factorial (108 scenarios and 540 simulations across
five seeds) when interaction effects are the target.

The $R=1$ setting was prespecified as a diagnostic boundary rather than a
confirmatory design.  Its three low--bulk--high samples and two regression
coefficients leave only one residual degree of freedom.  The corresponding
two-tail fit has two samples and two coefficients, leaves zero residual
degrees of freedom, and was therefore not fitted.  All
pairwise-design comparisons and the primary performance interpretation use
$R\geq3$; the one-replicate result is shown separately to expose the
consequence of insufficient replication.  BARCS-GC was evaluated only at
$R=1$ using Equations~\eqref{eq:guideconsistency}--\eqref{eq:gcz}; it was not
used to replace ordinary BARCS in supported replicated designs.

CRISPulator's arbitrary percentile ranges were set to $(0,0.25)$ for low,
$(0.75,1)$ for high, and $(0,1)$ for bulk.  Thus bulk deliberately overlaps
the tails and is not a third component of a multinomial partition.  An input
sample was generated as an independent sequencing draw from the same guide
frequency vector as bulk.  Low and high received phenotype scores equal to
the means of their corresponding standard-normal intervals,
\[
 d_{\mathrm{low}}=-1.271,\qquad
 d_{\mathrm{high}}=1.271,
\]
while bulk received $d_{\mathrm{bulk}}=0$.  BARCS and MAGeCK-MLE fitted this
score with replicate indicators \citep{li2015mageckvispr}.
Tail-only fits removed bulk without changing the low and high counts.  BARCS
guide tests were scaled with the simulated negative-control guides before
directional Stouffer aggregation; MAGeCK-MLE used its native gene-level Wald
result.  The separate 400-gene, five-seed baseline includes edgeR-QL, DESeq2,
and limma--voom in addition to BARCS and MAGeCK-MLE.  It is used to prevent the
restricted genome-scale set from supporting a general method-ranking claim.
CRISPhieRmix \citep{daley2018crisphiermix} was run on guide-level log$_2$
fold changes from a DESeq2 Wald fit of the same design
\citep{love2014deseq2}, which is the input its documentation specifies;
DESeq2 therefore appears as a preprocessing step in the 10,000-gene run, while
its native result is scored in the separate 400-gene comparison.
CRISPhieRmix was run with its bimodal alternative, because this screen
contains both phenotype-increasing and phenotype-decreasing genes and the
one-sided default would score only one direction.  Its gene effect is the
mean guide log$_2$ fold change and its local false-discovery rate stands in
for a $p$-value wherever a ranking is required, since the model reports
neither a $p$-value nor a signed gene effect.  Reported runtimes cover the
primary low--bulk--high model fit and exclude simulation, file input/output,
and guide-to-gene aggregation.

Genes from the increasing and decreasing classes were positives.  Ranking
used the two-sided gene significance score.  AUROC and average precision
measure active-gene discrimination.  Effect recovery is Spearman correlation
between the inferred gene effect and CRISPulator's median theoretical guide
phenotype among active genes.  Directional recall is the fraction of active
genes with gene FDR below 0.10 and an effect sign matching the simulated
class.  Realized FDP is the fraction of called genes that were inactive, and
F1 uses active versus inactive status at the same FDR threshold.
